\documentclass[11pt]{article}

\usepackage[margin=1in]{geometry}
\usepackage{amsmath,amssymb,amsthm,mathtools}
\usepackage{graphicx}
\usepackage[hidelinks]{hyperref}

\newtheorem{theorem}{Theorem}
\newtheorem{lemma}{Lemma}
\newtheorem{corollary}{Corollary}
\newtheorem{remark}{Remark}

\newcommand{\R}{\mathbb R}
\newcommand{\Loe}{\mathcal L}
\newcommand{\Pos}{(0,\infty)}

\title{Rational L\"owner Approximation of Positive Monomials}
\author{Automated Math Discovery Candidate Proof}
\date{August 9, 2026}

\begin{document}
\maketitle

\begin{center}
\fbox{\parbox{0.91\textwidth}{\textbf{Status.}
Candidate full solution, likely valid pending expert review.  We prove a
stronger multivariable theorem: every positive monomial of total exponent at
most one is a compact-open limit of rational globally operator-monotone
functions.}}
\end{center}

\section*{Source Question}

Agler, McCarthy, and Young~\cite{AMY} prove that
$(z^1z^2)^s$ is globally operator monotone on $\Pos^2$ for
$0\leq s\leq \tfrac12$.  Immediately after that example they ask:

\begin{quote}
\emph{We do not know if $(z^1z^2)^s$ can be approximated by rational
functions in the L\"owner class.}
\end{quote}

The question occurs on source page 49 (PDF page 49) in the final section of
arXiv:1009.3921v3, immediately after Example~9.8 begins on page 48; in the
source archive this is \texttt{glob.tex}, lines 247--290.

\begin{center}
\includegraphics[width=\textwidth]{figures/open_problem_context_crop.png}

\medskip
\includegraphics[width=\textwidth]{figures/open_problem_crop.png}
\end{center}

We use the standard compact-open meaning of approximation.  Thus a sequence
on $\Pos^d$ approximates a function if it converges uniformly on every compact
subset.

\section*{Definitions}

Let $E\subset\R^d$ be open.  A real function $f$ on $E$ is
\emph{globally operator monotone} if
\[
 S^r\leq T^r\quad(1\leq r\leq d)
 \qquad\Longrightarrow\qquad f(S)\leq f(T)
\]
whenever $S$ and $T$ are commuting $d$-tuples of bounded self-adjoint
operators with joint spectra in $E$.  The coordinates of $S$ commute with
one another, as do those of $T$; no commutation between the two tuples is
assumed.

The L\"owner class $\Loe(E)$ is the boundary-value class defined in
\cite{AMY}.  We need only their Theorem~8.1: for real $C^1$ functions on
$E$, membership in $\Loe(E)$ is equivalent to local operator monotonicity.
In particular,
\begin{equation}\label{eq:global-to-loewner}
 \text{globally operator monotone on }E
 \quad\Longrightarrow\quad f\in\Loe(E).
\end{equation}

\section*{Proof Intuition}

The source function looks nonrational because it contains fractional powers.
The useful observation is that fractional powers and weighted geometric means
are positive averages of very simple rational order-preserving maps.

For $0<\theta<1$, the scalar weighted geometric mean has the integral
\[
 x^{1-\theta}y^\theta
 =\frac{\sin(\pi\theta)}{\pi}
   \int_0^\infty \lambda^{\theta-1}
   \bigl(x^{-1}+\lambda y^{-1}\bigr)^{-1}\,d\lambda.
\]
Each integrand is jointly operator monotone: inversion reverses order, the
positive sum preserves it, and the final inversion reverses it again.
Similarly, $t^p$ is a positive integral of the rational operator-monotone
maps $t/(t+\lambda)$.  Finite positive quadrature sums therefore give
rational operator-monotone approximants.

Nested composition causes no loss.  Functions of a fixed commuting tuple
commute, and a globally operator-monotone binary map can be applied to any
ordered pair of such outputs.  This lets us build an arbitrary weighted
geometric monomial recursively and then apply the outer power.  The same
positive integral formulas give uniform convergence on compact sets.

\section*{Main Theorem}

\begin{theorem}\label{thm:main}
Let $d\geq1$ and let $a_1,\ldots,a_d\geq0$ satisfy
\[
                         p:=\sum_{j=1}^d a_j\leq1.
\]
There is a sequence of rational functions $R_n$ on $\R^d$, with no poles on
$\Pos^d$, such that:
\begin{enumerate}
\item each $R_n$ is globally operator monotone on $\Pos^d$;
\item each $R_n$ belongs to the L\"owner class $\Loe(\Pos^d)$;
\item $R_n\to\prod_{j=1}^d x_j^{a_j}$ uniformly on every compact subset of
$\Pos^d$.
\end{enumerate}
\end{theorem}

\begin{corollary}\label{cor:source}
For every $0\leq s\leq\tfrac12$, the function
$(z^1z^2)^s$ on $\Pos^2$ is a compact-open limit of rational functions in
$\Loe(\Pos^2)$.  Hence the question in \cite{AMY} has an affirmative answer.
\end{corollary}

\section*{Rational Building Blocks}

\begin{lemma}[Positive quadrature for a power]\label{lem:power}
For $0<p<1$, the function $t^p$ on $\Pos$ is a compact-open limit of finite
positive sums of the rational functions
\[
                         q_\lambda(t)=\frac{t}{t+\lambda},
                         \qquad \lambda>0.
\]
Every such finite positive sum is operator monotone and strictly positive on
$\Pos$.
\end{lemma}

\begin{proof}
Euler's beta integral gives
\begin{equation}\label{eq:power-integral}
 t^p=\frac{\sin(\pi p)}{\pi}
       \int_0^\infty \lambda^{p-1}\frac{t}{t+\lambda}\,d\lambda.
\end{equation}
Indeed, the substitution $\lambda=tu$ reduces the integral to
$t^p\int_0^\infty u^{p-1}(1+u)^{-1}du$.

The map $q_\lambda$ is operator monotone because
\[
 q_\lambda(A)=I-\lambda(A+\lambda I)^{-1},
\]
and inversion reverses the order of positive invertible operators.  Hence
finite positive sums of the $q_\lambda$ are operator monotone.

For $t\in[c,C]\subset\Pos$, the part of \eqref{eq:power-integral} over
$(0,\delta)$ is at most $\delta^p/p$, uniformly in $t$, and the part over
$(M,\infty)$ is at most $C M^{p-1}/(1-p)$.  On $[\delta,M]$ ordinary
positive Riemann sums converge uniformly for $t\in[c,C]$.  Letting
$\delta\downarrow0$, $M\uparrow\infty$, and the mesh tend to zero produces a
single sequence converging on every compact subinterval of $\Pos$.
\end{proof}

\begin{lemma}[Positive quadrature for a weighted geometric mean]
\label{lem:geometric}
Let $0<\theta<1$.  The function
\[
                         g_\theta(x,y)=x^{1-\theta}y^\theta
\]
on $\Pos^2$ is a compact-open limit of finite positive sums of the rational
functions
\[
 h_\lambda(x,y)
   =\bigl(x^{-1}+\lambda y^{-1}\bigr)^{-1}
   =\frac{xy}{y+\lambda x},
 \qquad \lambda>0.
\]
Every such sum is globally operator monotone and strictly positive on
$\Pos^2$.
\end{lemma}

\begin{proof}
The beta integral and the substitution $\lambda=(y/x)u$ give
\begin{equation}\label{eq:geometric-integral}
 x^{1-\theta}y^\theta
 =\frac{\sin(\pi\theta)}{\pi}
   \int_0^\infty \lambda^{\theta-1}
        \frac{xy}{y+\lambda x}\,d\lambda.
\end{equation}

Suppose $A_1\leq B_1$ and $A_2\leq B_2$ are positive invertible operators.
Then
\[
 A_1^{-1}+\lambda A_2^{-1}
 \geq B_1^{-1}+\lambda B_2^{-1},
\]
so another inversion yields
\[
 (A_1^{-1}+\lambda A_2^{-1})^{-1}
 \leq(B_1^{-1}+\lambda B_2^{-1})^{-1}.
\]
When each pair is commuting, these are exactly the joint functional-calculus
values of $h_\lambda$.  Thus $h_\lambda$ is globally operator monotone.

For $(x,y)\in[c,C]^2$, the integrand without the constant in
\eqref{eq:geometric-integral} is bounded by
$C\lambda^{\theta-1}$ near zero and by
$C\lambda^{\theta-2}$ near infinity.  The tails therefore tend uniformly to
zero.  Positive Riemann sums on a truncated compact $\lambda$-interval give
the asserted rational approximants and compact-open convergence.
\end{proof}

\section*{Composition and Proof of the Theorem}

\begin{lemma}[Order-preserving composition]\label{lem:composition}
Let $U,V$ be strictly positive globally operator-monotone functions on
$\Pos^d$, and let $H$ be globally operator monotone on $\Pos^2$.  Then
\[
                         x\longmapsto H(U(x),V(x))
\]
is globally operator monotone on $\Pos^d$.  If all three functions are
rational and have no poles on their positive domains, then so is the
composition.
\end{lemma}

\begin{proof}
For ordered commuting $d$-tuples $S\leq T$, global monotonicity gives
$U(S)\leq U(T)$ and $V(S)\leq V(T)$.  The pair $(U(S),V(S))$ commutes because
both entries belong to the joint functional calculus of $S$; the analogous
pair for $T$ also commutes.  The defining property of $H$ now gives
\[
 H(U(S),V(S))\leq H(U(T),V(T)).
\]
The rationality assertion is immediate from substitution, and positivity
keeps all positive-domain denominators away from zero.
\end{proof}

\begin{proof}[Proof of Theorem~\ref{thm:main}]
The case $p=0$ is the constant function $1$.  Coordinates with $a_j=0$ may
be omitted, so assume first that the active exponents are positive.  Put
\[
                         w_j=\frac{a_j}{p},
                         \qquad \sum_j w_j=1.
\]
We first approximate the normalized weighted monomial
$G(x)=\prod_jx_j^{w_j}$.

Order the active coordinates and set $W_k=\sum_{j=1}^k w_j$.  Recursively
define
\[
 G_1(x)=x_1,
 \qquad
 G_k(x)=g_{w_k/W_k}\bigl(G_{k-1}(x),x_k\bigr)
 \quad (k\geq2).
\]
A scalar exponent calculation gives
\[
                         G_k(x)=\prod_{j=1}^k x_j^{w_j/W_k};
\]
hence $G_d=G$.

For every binary $g_{w_k/W_k}$ choose the rational positive quadrature
approximants from Lemma~\ref{lem:geometric}, and substitute them recursively.
Lemma~\ref{lem:composition} shows that every resulting nested rational
function is globally operator monotone.  A standard diagonal choice makes
these nested functions converge to $G$ uniformly on compact subsets.  For
completeness, the induction is harmless: on any fixed compact domain,
$G_{k-1}$ has a compact range bounded away from zero; uniform convergence of
the preceding stage puts all sufficiently late approximating ranges in a
slightly larger compact subinterval of $\Pos$, where the next binary
quadrature converges uniformly.

If $p=1$, these nested approximants already converge to
$G=\prod_jx_j^{a_j}$.  If $0<p<1$, compose them with the positive rational
operator-monotone approximants to $t^p$ from Lemma~\ref{lem:power}.  Another
application of Lemma~\ref{lem:composition}, followed by the same compact-range
diagonal argument, yields rational globally operator-monotone functions
converging compact-open to
\[
                         G(x)^p=\prod_{j=1}^d x_j^{a_j}.
\]

Finally, every approximant belongs to $\Loe(\Pos^d)$ by
\eqref{eq:global-to-loewner}.  This proves all three assertions.
\end{proof}

\begin{proof}[Proof of Corollary~\ref{cor:source}]
Take $d=2$ and $a_1=a_2=s$.  The hypothesis of
Theorem~\ref{thm:main} is exactly $2s\leq1$.  The endpoint $s=0$ is the
constant $1$, so the entire range in the source question is covered.
\end{proof}

\section*{Verification and Limitations}

No numerical or symbolic computation is used.  The proof reduces to the two
positive integral formulas \eqref{eq:power-integral} and
\eqref{eq:geometric-integral}, two order reversals under inversion, and a
composition argument inside joint functional calculus.

The packet proves locally uniform approximation on $\Pos^d$.  It does not
claim uniform approximation up to the boundary, where the monomials may lose
smoothness, or a degree-optimal approximation rate.

A bounded search on August 9, 2026 used the exact source phrase and close
variants involving rational L\"owner approximation, operator-monotone
monomials, and weighted geometric/harmonic means.  It found the source paper
and later one-variable approximation work, but no explicit answer to the
question or statement of Theorem~\ref{thm:main}.  Since all ingredients are
classical, novelty confidence is moderate-to-low pending specialist review.

\begin{thebibliography}{9}

\bibitem{AMY}
J.~Agler, J.~E. McCarthy, and N.~J. Young,
\emph{Operator monotone functions and L\"owner functions of several
variables}, Ann. of Math. (2) \textbf{176} (2012), 1783--1826;
arXiv:1009.3921.

\end{thebibliography}

\end{document}
